\bookchapter{Глава 31: Съёмка коротких видео и введение в создание контента}

\chapterintro{В 2026 году короткие видео уже не «игрушка молодёжи» — это базовый способ получения информации и самовыражения для всех. Хотите ли вы фиксировать жизнь, делиться знаниями или подрабатывать — телефон станет лучшим инструментом для творчества: не нужно покупать камеру, не нужно осваивать Premiere, достаточно одного телефона. Эта глава поможет начать с нуля.}

\booksection{I. Прежде чем снимать, решите: для кого вы снимаете?}

Первая ошибка многих при съёмке видео — «взял телефон и снимаю», а потом непонятно, для кого. Прежде чем начать, ответьте на три вопроса:

\begin{enumerate}
  \item \textbf{Кто моя аудитория?}— Снимать для детей, для ровесников или для старшего поколения — это совершенно разные стиль и содержание.
  \item \textbf{Какую ценность я даю?}— Чтобы рассмешить (развлечение), чтобы научить (знания), чтобы вызвать аппетит (еда), чтобы тронуть (эмоции). Должно быть хотя бы одно.
  \item \textbf{На какой платформе я смогу регулярно обновляться?}— Douyin/Kuaishou (общее развлечение), Xiaohongshu (образ жизни), Bilibili (глубокий контент), Shipinhao (социальный круг). Выберите одну основную площадку и развивайтесь на ней, не пытайтесь сразу выходить на все платформы.
\end{enumerate}

\booksection{II. Базовые настройки съёмки на телефоне — 5 обязательных действий перед съёмкой}

\begin{enumerate}
  \item \textbf{Разрешение 1080P, частота кадров 30 fps}: Это стандарт платформ коротких видео. 4K-файлы слишком большие, загрузка медленная, а 60 fps на социальных платформах не даёт преимущества, лишь увеличивает файл.
  \item \textbf{Протрите объектив}: Первый и самый важный пункт. Перед съёмкой протрите объектив тканью — картинка сразу станет чётче.
  \item \textbf{Заблокируйте фокус и экспозицию}: Долгим нажатием на объект в видоискателе в течение 2 секунд появится подсказка «Блокировка AE/AF». Так при движении картинка не будет то яркой, то тёмной.
  \item \textbf{Горизонтально или вертикально?}— Для Douyin/Kuaishou/Shipinhao/Xiaohongshu снимайте\textbf{вертикально}(9:16), для Bilibili/YouTube —\textbf{горизонтально}(16:9). Сначала выберите платформу, потом ориентацию.
  \item \textbf{Включите сетку}: Настройки > Камера > Сетка. Это основа композиции по правилу третей, помогает быстро найти центр кадра.
\end{enumerate}

\booksection{III. Пять приёмов движения камеры — оживите кадр}

Главная проблема новичков — слишком «статичный» кадр: всю съёмку стоят на месте. Освойте пять базовых движений камеры, и картинка сразу станет профессиональной:

\begin{itemize}
  \item \textbf{Налезд (Push)}: Телефон медленно приближается к объекту съёмки. Создаёт чувство «приближения» — напряжения или интимности.
  \item \textbf{Отъезд (Pull)}: Телефон медленно удаляется от объекта. От детали к общему плану, показывает обстановку.
  \item \textbf{Панорама (Pan)}: Телефон стоит на месте, поворачивается влево-вправо. Подходит для съёмки пейзажей и показа пространства.
  \item \textbf{Проехать (Track)}: Телефон движется параллельно, следуя за объектом. Нужна уверенная рука, лучше со стабилизатором.
  \item \textbf{Слежение (Follow)}: Телефон идёт позади или впереди объекта. Самый частый приём в Vlog.
\end{itemize}

\textbf{Плавность — главное}: Дрожащая картинка выглядит непрофессионально. Двумя руками держать телефон надёжнее, чем одной; при ходьбе слегка согните колени и используйте «походку ниндзя», чтобы уменьшить тряску; если позволяет бюджет, купите стабилизатор для смартфона (серия DJI OM — около 500–800 юаней) — качество кардинально улучшится.

\booksection{IV. Введение в монтаж — достаточно CapCut}

CapCut (китайская версия Jianying) — лучший инструмент для монтажа на телефоне для новичков: бесплатный, с полным функционалом и множеством шаблонов. Базовый рабочий процесс:

\begin{enumerate}
  \item \textbf{Импорт материала}: Откройте CapCut > «Начать проект» > выберите снятые видео и импортируйте их.
  \item \textbf{Обрезка}: Перетаскивайте таймлайн, обрезайте начало и конец каждого фрагмента. Принцип: оставляйте только содержательные части, лучше короче, чем длиннее.
  \item \textbf{Добавьте переходы}: Нажмите кнопку «Переход» между двумя клипами и выберите простой эффект. Новичкам рекомендуется «наплыв» или «затемнение», не используйте броские 3D-переходы.
  \item \textbf{Добавьте субтитры}: В CapCut есть функция «распознавание субтитров» — одним нажатием автоматически генерирует субтитры. Точность около 90\%, достаточно проверить вручную.
  \item \textbf{Добавьте музыку}: В музыкальной библиотеке CapCut много музыки, свободной от авторских прав. Выберите BGM, ритм которого соответствует содержанию, громкость установите примерно на 1/3 от оригинального звука видео.
  \item \textbf{Экспорт}: 1080P, частота 30 fps, битрейт — «рекомендуемый».
\end{enumerate}

\booksection{V. Важность обложки — 70\% людей смотрят только по обложке}

Обложка короткого видео — это «витрина», она решает, зайдёт ли зритель. Несколько советов по обложке:

\begin{itemize}
  \item \textbf{На обложке — одна фраза}: Крупно укажите суть видео, например «Освойте обработку фото на телефоне за 3 минуты». Текст должен быть крупным, ярким, контрастировать с фоном.
  \item \textbf{Обложка должна соответствовать содержанию}: Не будьте «кликбейтом» — обложка одно, а внутри другое, зрители тут же уйдут, а алгоритм снизит рейтинг.
  \item \textbf{Выразительные эмоции}: Если на обложке есть человек, выберите момент с живой или утрированной мимикой. Лицо привлекает клики лучше текста.
  \item \textbf{Делайте обложку в CapCut}: Выберите лучший кадр в CapCut, добавьте наложение текста и экспортируйте как изображение-обложку.
\end{itemize}

\booksection{VI. Три принципа создания контента}

\begin{itemize}
  \item \textbf{Первые 3 секунды решают всё}: На платформах коротких видео зрители решают за 3 секунды, уйти или продолжить смотреть. В первые 3 секунды нужна «крючок» — вопрос, интрига или неожиданный кадр.
  \item \textbf{Одно видео — одна тема}: Самая частая ошибка новичков — запихнуть в одно видео слишком много. Одно видео должно решать одну задачу; раскрыть тему глубоко важнее, чем охватить всё.
  \item \textbf{Постоянство важнее таланта}: Алгоритму рекомендаций нужно «раскачать аккаунт». Первые 20 видео почти не наберут просмотров — это нормально. Продолжайте выкладывать 2–3 видео в неделю, чтобы алгоритм понял направление вашего контента, и трафик постепенно вырастет.
\end{itemize}

\begin{tcbpractice}{Практика: снимите 30-секундное короткое видео}
\begin{enumerate}
  \item Выберите небольшое знание, в котором разбираетесь (например, «как увеличить шрифт на телефоне»).
  \item Снимите вертикально на телефон 3–5 фрагментов по 5–10 секунд каждый.
  \item В CapCut склейте, обрежьте, добавьте субтитры и BGM.
  \item Сделайте обложку с крупным текстом.
  \item После экспорта сначала покажите другу, соберите отзывы, а потом публикуйте.
\end{enumerate}

\end{tcbpractice}

\begin{tcbsummary}{Итоги главы}
\begin{itemize}
  \item Перед съёмкой решите: кто аудитория, какую ценность даёте, на какой платформе публикуете
  \item 5 шагов перед съёмкой: 1080P/30 fps, протрите объектив, заблокируйте фокус и экспозицию, выберите ориентацию, включите сетку
  \item Пять приёмов движения камеры: наезд, отъезд, панорама, проезд, слежение — плавность прежде всего
  \item Монтаж в CapCut: обрезка переход субтитры BGM экспорт, шесть шагов
  \item Первые 3 секунды решают всё, одно видео — одна тема, постоянство важнее таланта
\end{itemize}

\end{tcbsummary}
